\subsection{Semantic Regularities in Embedding Spaces}\label{sec:semantic-regularities-in-embedding-spaces}

\subsubsection{Semantic Relationships}\label{sec:semantic-relationships}

After training Word2Vec on a sufficiently large corpus, the learned embedding space exhibits an unexpected and remarkable property: \hl{semantic relationships between words become encoded as vector directions.}

\highspace
Recall that each word is represented by a dense vector:
\begin{equation*}
    \mathbf{w} \in \mathbb{R}^{m}
\end{equation*}
Where nearby vectors correspond to words that tend to appear in similar contexts.

\highspace
Initially, researchers expected only that semantically similar words would form clusters, for example:
\begin{itemize}
    \item \emph{dog}, \emph{wolf}, \emph{cat};
    \item \emph{France}, \emph{Germany}, \emph{Italy};
    \item \emph{eat}, \emph{drink}, \emph{cook}.
\end{itemize}
However, Word2Vec revealed something much stronger: \hl{many semantic relationships correspond to approximately constant vector offsets throughout the embedding space.} In other words, not only are similar words close together, but the \textbf{difference vector} between two related words is often nearly the same regardless of the specific words involved. This phenomenon is referred to as a \textbf{semantic regularity} in the embedding space.

\highspace
\begin{flushleft}
    \textcolor{Green3}{\faIcon{square-root-alt} \textbf{Geometric Interpretation}}
\end{flushleft}
Suppose each word is represented as a point in a high-dimensional vector space. Instead of looking only at the positions of the points, we can also consider the \textbf{vector connecting two words}. For two words $a$ and $b$, this relationship is represented by the vector:
\begin{equation*}
    \overrightarrow{ab} = \mathbf{w}_b - \mathbf{w}_a
\end{equation*}
If two different pairs of words express the same semantic relationship, their difference vectors tend to be almost parallel and have similar magnitudes. This means that semantic \hl{information is encoded} not only in the \hl{location of individual words} but also in the \hl{\textbf{directions} connecting them}.

\highspace
\begin{flushleft}
    \textcolor{Green3}{\faIcon{search} \textbf{Country-Capital Relationships}}
\end{flushleft}
One of the first observations made by Mikolov and colleagues concerns country-capital pairs. For example:
\begin{align*}
    \mathbf{w}_{Paris} - \mathbf{w}_{France} &\approx
    \mathbf{w}_{Rome} - \mathbf{w}_{Italy} \\
    &\approx \mathbf{w}_{Berlin} - \mathbf{w}_{Germany} \\
    &\approx \mathbf{w}_{Madrid} - \mathbf{w}_{Spain} \\
    &\approx \mathbf{w}_{Lisbon} - \mathbf{w}_{Portugal}
\end{align*}
Although these countries occupy \textbf{different regions of the embedding space}, the vectors connecting each country to its capital point in \textbf{approximately the same direction}.

\begin{figure}[!htp]
    \centering
    \resizebox{0.78\textwidth}{!}{%
    \begin{tikzpicture}[
        >=Latex,
        font=\small,
        every node/.style={align=center},
        country/.style={
            circle,
            fill=blue!70,
            inner sep=2.4pt
        },
        capital/.style={
            circle,
            fill=orange!85!black,
            inner sep=2.4pt
        },
        relation/.style={
            ->,
            very thick,
            blue!60!black
        },
        guide/.style={
            ->,
            thick,
            dashed,
            gray!70
        },
        axis/.style={
            ->,
            gray!70
        }
    ]

    % ============================================================
    % AXES
    % ============================================================
    \draw[axis] (-0.3,0) -- (7.4,0) node[right] {$x_1$};
    \draw[axis] (0,-0.3) -- (0,4.8) node[above] {$x_2$};

    % ============================================================
    % POINTS: COUNTRIES
    % ============================================================
    \node[country, label=left:{France}]   (france)   at (1.0,0.9) {};
    \node[country, label=left:{Italy}]    (italy)    at (1.4,1.9) {};
    \node[country, label=left:{Germany}]  (germany)  at (1.9,3.0) {};
    \node[country, label=left:{Spain}]    (spain)    at (2.5,0.45) {};
    \node[country, label=left:{Portugal}] (portugal) at (2.9,1.35) {};

    % ============================================================
    % POINTS: CAPITALS
    % ============================================================
    \node[capital, label=right:{Paris}]   (paris)   at (4.4,1.45) {};
    \node[capital, label=right:{Rome}]    (rome)    at (4.8,2.45) {};
    \node[capital, label=right:{Berlin}]  (berlin)  at (5.3,3.55) {};
    \node[capital, label=right:{Madrid}]  (madrid)  at (5.9,1.00) {};
    \node[capital, label=right:{Lisbon}]  (lisbon)  at (6.3,1.90) {};

    % ============================================================
    % RELATION VECTORS: COUNTRY -> CAPITAL
    % ============================================================
    \draw[relation] (france) -- (paris);
    \draw[relation] (italy) -- (rome);
    \draw[relation] (germany) -- (berlin);
    \draw[relation] (spain) -- (madrid);
    \draw[relation] (portugal) -- (lisbon);

    % ============================================================
    % COMMON OFFSET VECTOR
    % ============================================================
    \draw[guide] (0.8,4.35) -- ++(2.05,0.35)
        node[midway, above, font=\scriptsize]
        {$\Delta_{\text{capital}}$};

    \node[
        font=\scriptsize,
        text width=3.8cm,
        align=left
    ] at (4.85,4.35) {
        approximately parallel offsets
    };

    % ============================================================
    % LEGEND
    % ============================================================
    \node[
        draw=gray!40,
        rounded corners,
        fill=gray!3,
        inner xsep=8pt,
        inner ysep=6pt
    ] at (3.7,-0.85) {
        \begin{tabular}{@{}cl@{\qquad}cl@{}}
            \tikz{\node[country] {}; } & country
            &
            \tikz{\node[capital] {}; } & capital
        \end{tabular}
    };

    \end{tikzpicture}%
    }
    \caption{Geometric interpretation of semantic regularities in Word2Vec. Country--capital relationships are represented by approximately parallel vector offsets in the embedding space.}
    \label{fig:word2vec-country-capital-regularity}
\end{figure}

\noindent
\textcolor{Green3}{\faIcon{question-circle} \textbf{Why does this happen?}} The network is \textbf{never explicitly told} that Paris is the capital of France or that Rome is the capital of Italy. Instead, it repeatedly observes contexts such as:
\begin{itemize}
    \item \emph{Paris is the capital of France.}
    \item \emph{Rome is the capital of Italy.}
    \item \emph{Madrid is the capital of Spain.}
\end{itemize}
As these contextual patterns appear throughout the corpus, gradient updates gradually arrange the embeddings so that the same \textbf{semantic relationship is represented by similar vector offsets}. Consequently, the notion of ``\emph{capital of}'' emerges automatically from the training data.

\highspace
\textcolor{Green3}{\faIcon{link} \textbf{King-Queen Relationships.}} Another famous example concerns \textbf{gender relationships}. The vector:
\begin{equation*}
    \mathbf{w}_{\text{king}} - \mathbf{w}_{\text{queen}}
\end{equation*}
Is approximately parallel to:
\begin{equation*}
    \mathbf{w}_{\text{prince}} - \mathbf{w}_{\text{princess}}
\end{equation*}
And similarly to:
\begin{equation*}
    \mathbf{w}_{\text{man}} - \mathbf{w}_{\text{woman}}
\end{equation*}
Has nearly the same direction. Again, Word2Vec has not been taught what ``\emph{gender}'' means. It only learns from contextual information, yet this semantic concept naturally emerges as a consistent direction in the embedding space.

\highspace
\begin{flushleft}
    \textcolor{Green3}{\faIcon{question-circle} \textbf{Why are these relationships important?}}
\end{flushleft}
These semantic regularities demonstrate that Word2Vec learns \textbf{more than word similarity}. Earlier embedding methods primarily captured the idea that \emph{similar words are close together}. Word2Vec goes one step further: \hl{relationships between words become linear transformations in the embedding space.} This discovery showed that semantic knowledge can be represented geometrically, allowing simple vector operations to reveal complex linguistic relationships. This observation directly motivates the next section, where we will see that these vector relationships make it possible to solve analogies such as:
\begin{equation*}
    \mathbf{w}_{\text{king}} - \mathbf{w}_{\text{man}} + \mathbf{w}_{\text{woman}} \approx \mathbf{w}_{\text{queen}}
\end{equation*}
Using nothing more than vector arithmetic.